\section*{Введение}
\addcontentsline{toc}{section}{Введение}

В главе~2 мы научились строить поле разложения многочлена и считать гомоморфизмы между такими полями. Теперь главным
объектом становится \emph{группа} $F$-автоморфизмов поля $E$. Основная теорема теории Галуа утверждает, что для
конечного расширения Галуа $E/F$ (например, поля разложения сепарабельного многочлена) промежуточные поля $F\subset M\subset E$
взаимно однозначно соответствуют подгруппам группы $\Gal(E/F)$: вопросы о полях превращаются в вопросы о конечных
группах. Инструменты прежние~--- подсчёт гомоморфизмов (\rf{2.12}, \rf{2.14}) и линейная алгебра; новый ключевой факт~---
теорема Артина о неподвижном поле конечной группы.

\begin{itemize}[leftmargin=*]
\item \emph{Группы автоморфизмов полей} (\rf{3.1}--\rf{3.5c}). Примеры: $\Aut(\CC/\R)$, $\Aut(\R)=1$ (\rf{3.1d}), группа
$\Aut(k(X)/k)\cong PGL_2(k)$ дробно-линейных подстановок (\rf{3.1b}) с вычислением степени $[k(X):k(u)]$ для рациональной
функции $u$ (\rf{3.1a}), сфера Римана и группа Кремоны (\rf{3.1c}). Для поля разложения сепарабельного многочлена
$|\Aut(E/F)|=[E:F]$ (\rf{3.2}), и примеры, показывающие, что оба условия нужны (\rf{3.3}); число автоморфизмов простого расширения (\rf{3.3a}). Неподвижное поле $E^G$ и
формальные свойства соответствия <<подгруппа $\leftrightarrow$ подполе>> (\rf{3.3b}, \rf{3.3c}). Теорема Артина
$[E:E^G]\le(G:1)$ (\rf{3.4}) и её следствие $G=\Aut(E/E^G)$ для конечной группы $G$ (\rf{3.5}); почему это не тавтология
(\rf{3.5a}) и почему для бесконечных групп утверждение неверно (\rf{3.5b}); нахождение неподвижных полей в $k(X)$, в том
числе для группы порядка $6$, переставляющей $0,1,\infty$ (\rf{3.5c}).
\item \emph{Сепарабельные, нормальные расширения и расширения Галуа} (\rf{3.6}--\rf{3.16}). Определения и примеры
($\Q[\sqrt[3]2]/\Q$ не нормально, $\F_p(T)/\F_p(T^p)$ не сепарабельно, \rf{3.8}). Четыре равносильных определения расширения
Галуа (\rf{3.10}); конечные нормальные расширения~--- это в точности поля разложения многочленов (\rf{3.10a}); минимальный
многочлен элемента $h(\alpha)$ через корни $\mu_\alpha$ (\rf{3.10b}); теорема Артина в окончательной форме (\rf{3.11}); вложение
конечного сепарабельного расширения в расширение Галуа (\rf{3.12}) и нормальное замыкание (\rf{3.12a}); поведение при
промежуточных полях (\rf{3.13}) и нетранзитивность нормальности (\rf{3.13a}). Сопряжённые элементы (\rf{3.14}, \rf{3.14a}).
Сепарабельное замыкание, сепарабельная степень, чисто несепарабельные расширения (\rf{3.15}, \rf{3.15a}); критерии
сепарабельности (\rf{3.15b}), транзитивность сепарабельности с явной формулой для минимального многочлена над меньшим
полем (\rf{3.15c}, \rf{3.15d}), связь с этальными алгебрами (\rf{3.15e}). Циклические, абелевы и разрешимые расширения
(\rf{3.16}).
\item \emph{Основная теорема теории Галуа} (\rf{3.17}--\rf{3.21c}). Соответствие Галуа $H\mapsto E^H$, $M\mapsto\Gal(E/M)$ и его
свойства: обращение включений, индексы и степени, сопряжённые подгруппы и сопряжённые поля, нормальные подгруппы и
нормальные подрасширения (\rf{3.17}, \rf{3.17a}); $F$-вложения подрасширения в $E$ и смежные классы (\rf{3.17b}); почему нужно
условие <<расширение Галуа>> (\rf{3.17c}). Композиты и пересечения, нормальное замыкание на языке подгрупп (\rf{3.18}, \rf{3.18a});
поднятие расширения Галуа на композит (\rf{3.19}, \rf{3.19a}) и степень композита (\rf{3.20}; без условия Галуа~--- лишь неравенство, \rf{3.20a}); группа Галуа композита двух расширений
Галуа как расслоённое произведение групп (\rf{3.21}), прямое произведение при $E_1\cap E_2=F$ (\rf{3.21a}) и связь с
кодекартовыми квадратами полей (\rf{3.21b}; ср.\ \rf{2.15a}--\rf{2.15b}), второе доказательство \rf{3.21} и склейки поля Галуа с самим собой (\rf{3.21c}).
\item \emph{Примеры} (\rf{3.22}--\rf{3.23e}). Круговое поле $\Q[\zeta_7]$: решётка подполей, гауссовы периоды и квадратичное
подполе $\Q[\sqrt{-7}]$, и то же для $\Q[\zeta_p]$ в общем виде (\rf{3.22}; квадратичное подполе подробно~--- в упражнении
3-2); кубическое поле $\Q[\cos\frac{2\pi}9]$ и многочлен $X^3-3X+1$ (\rf{3.22a}). Вычисление гауссовых периодов подсчётом
(\rf{3.22b}): кубическое подполе поля $\Q[\zeta_{13}]$ (\rf{3.22c}) и его подполе степени $4$, кубические периоды для простых
$p\equiv1\pmod3$; формула Кардано для периодов (\rf{3.22d}) и неприводимый
случай: корни неприводимого кубического многочлена с тремя вещественными корнями не выражаются через вещественные
радикалы (\rf{3.22e}, \rf{3.22f}). Поле разложения $X^5-2$ над $\Q$, его группа порядка $20$~--- группа Фробениуса
$\Z/5\rtimes\Z/4$ аффинных преобразований прямой $\F_5$ (\rf{3.23}, \rf{3.23a})~--- и все его подполя (\rf{3.23b}). Подполя поля
разложения $X^3+3X+1$ (\rf{3.23c}). Многочлены Артина--Шрайера $X^p-X-a$: циклические расширения степени $p$, в том
числе над $\F_p$ и над $k(t)$ (\rf{3.23d}), и почему корни $X^p-X-t$ не лежат ни в какой радикальной башне (\rf{3.23e}).
\item \emph{Построимые числа ещё раз} (\rf{3.24}--\rf{3.27}). Число, лежащее в вещественном расширении Галуа поля $\Q$
степени $2^r$, построимо (\rf{3.24}); точный критерий: комплексное число $z$ построимо тогда и только тогда, когда поле
разложения его минимального многочлена имеет над $\Q$ степень $2^r$ (\rf{3.25b}), и почему вещественности требовать
нельзя (\rf{3.25c}). Правильный $p$-угольник для простых Ферма $p=2^k+1$ (\rf{3.26}) и явная башня квадратичных расширений
для семнадцатиугольника (\rf{3.26a}); необходимое условие <<степень~--- степень двойки>> не достаточно (\rf{3.27}).
\item \emph{Группа Галуа многочлена} (\rf{3.27a}--\rf{3.27h}). Группа $G_f$ сепарабельного многочлена как группа перестановок
его корней, сохраняющих все полиномиальные соотношения между ними (\rf{3.27b}); её порядок делит $(\deg f)!$, а орбиты~---
множества корней неприводимых множителей (\rf{3.27d}). Приводимые многочлены $f=f_1\cdots f_k$ (например, $f=gh$):
вложение $G_f\subset G_{f_1}\times\dots\times G_{f_k}$, его образ при $k=2$~--- расслоённое произведение (\rf{3.27e}, через
\rf{3.21}), почему не всякий набор поднимается (\rf{3.27f}), примеры (\rf{3.27g}) и разрешимость $G_f$ (\rf{3.27h}).
\item \emph{Разрешимость уравнений} (\rf{3.27i}--\rf{3.30d}). Теорема Галуа: в характеристике $0$ уравнение $f(X)=0$
разрешимо в радикалах тогда и только тогда, когда группа $G_f$ разрешима (\rf{3.28}; доказательство~--- в \rf{5.34});
формула Кардано для кубических уравнений (\rf{3.28a}). Что меняется в характеристике $p$ (\rf{3.29}; ср.\
\rf{5.35a}--\rf{5.35e}): корни многочлена Артина--Шрайера $X^p-X-t$ над $\Fpb(t)$ не лежат ни в какой радикальной башне,
хотя его группа Галуа~--- $\Z/p\Z$ (\rf{3.29a}). Уравнения простой степени: полное доказательство того, что транзитивная группа перестановок
простой степени $p$ разрешима тогда и только тогда, когда при подходящей нумерации она содержится в группе аффинных
преобразований $x\mapsto ux+v$ прямой $\F_p$, и тогда и только тогда, когда выполнено условие Абеля (\rf{3.30a}); отсюда неприводимое уравнение простой степени разрешимо в радикалах
тогда и только тогда, когда его поле разложения порождено любыми двумя корнями (\rf{3.30b}), и следствие Кронекера о
вещественных корнях (\rf{3.30c}). Плоскость Фано показывает, почему в условии Абеля нужны \emph{любые} три корня
(\rf{3.30d}).
\item \emph{Упражнения} 3-1--3-5 с решениями: $F(X^2)\cap F(X^2-X)=F$ (через группу, порождённую двумя инволюциями
$X\mapsto-X$ и $X\mapsto1-X$, и лемму \rf{3.1a}); квадратичное подполе $\Q[\zeta_p]$ и гауссовы суммы; кватернионная группа
как группа Галуа поля $\Q[\sqrt2,\sqrt3,\sqrt{(2+\sqrt2)(3+\sqrt3)}]$ (с леммой о том, когда башня расширений Галуа сама
является расширением Галуа, \rf{3.30e}); группа автоморфизмов $\Aut(E^H/F)\cong N_G(H)/H$; порядок $\Aut(E/F)$ делит $[E:F]$.
\end{itemize}

Всюду в этой главе $F$~--- поле, а $\mu_\alpha$ (или $\mu^F_\alpha$) обозначает минимальный многочлен элемента $\alpha$ над $F$.
